\documentclass{article}

\usepackage[final]{neurips_2023}

% suppress the "37th Conference on NeurIPS" footer
\makeatletter
\renewcommand{\@notice}{}
\makeatother

\usepackage[utf8]{inputenc}
\usepackage[T1]{fontenc}
\usepackage{hyperref}
\usepackage{url}
\usepackage{booktabs}
\usepackage{amsfonts}
\usepackage{amsmath}
\usepackage{amssymb}
\usepackage{nicefrac}
\usepackage{microtype}
\usepackage{xcolor}
\usepackage{graphicx}
\usepackage{multirow}
\usepackage{array}

\title{Graphical Model Structure Learning\\for Sparse Feature Circuits in GPT-2}

\author{%
  Siddharth Raj \\
  Toyota Technological Institute at Chicago \\
  \texttt{sidraj@uchicago.edu}
}

\begin{document}

\maketitle

% ─────────────────────────────────────────────────────────────────────────────
\begin{abstract}
% ─────────────────────────────────────────────────────────────────────────────
We ask whether probabilistic graphical model (PGM) structure learning,
applied to sparse feature activations, can automatically recover the circuits
inside a language model.
Using Sparse Autoencoders (SAEs) to extract 76 interpretable features from
GPT-2 Small's residual streams at layers 1 and 2, we treat feature activations
across 5,000 prompts as a multivariate dataset and run three algorithms: the
PC algorithm with Fisher-Z and G-test independence tests, a custom GES-BIC
implementation accelerated via Schur-complement score evaluations, and Mean
Ablation Patching as a causal proxy baseline.
We make two methodological contributions.
First, we introduce a dual-ground-truth evaluation protocol that separates
\emph{interventional} validation (zero-ablation) from \emph{observational}
validation (Spearman rank correlation), exposing a structural bias that
inflates Mean Ablation's apparent advantage.
Second, we show that GES-BIC dominates both PC variants on precision and
speed simultaneously: 157 edges, 0.667 ablation IA, in 36 seconds.
Pairwise Jaccard similarity of $\approx 0.10$ between PGM and ablation methods
demonstrates that statistical dependence and interventional sensitivity
discover fundamentally different edge sets.
Head-level aggregation reveals a consistent broad semantic circuit
(L0H9 $\to$ L1H2/H10) across all methods, but not the known induction circuit,
which we attribute to the feature selection regime rather than a failure of
structure learning.
\end{abstract}

% ─────────────────────────────────────────────────────────────────────────────
\section{Introduction}
% ─────────────────────────────────────────────────────────────────────────────

Understanding \emph{how} large language models implement their computations is
a central open problem. \emph{Mechanistic interpretability} decomposes model
behaviour into human-understandable circuits---subgraphs of components that
collectively perform a specific function \citep{olah2020zoom,
elhage2021framework}. Most circuit discovery to date has required substantial
manual effort: identifying components by hand, writing targeted prompts, and
iterating through ablation experiments \citep{wang2022ioi, olsson2022induction}.

Automating this process is the core motivation of our work. PGM structure
learning offers a principled framework: collect feature activations across
diverse inputs, model them as samples from a joint distribution over a DAG,
and recover the graph. If the model's circuits drive the statistical
dependencies in the activation data, structure learning should recover them.
Two challenges complicate this picture. First, circuit-specific features may
fire rarely, making them invisible to density-based feature filters. Second,
the standard evaluation metric---intervention agreement---conflates methods
that find statistical co-occurrence with those that directly measure causal
sensitivity.

We address both challenges. We apply PC (Fisher-Z and G-test) and GES-BIC to
SAE feature activations from GPT-2 Small, with Mean Ablation Patching as a
causal baseline. We introduce a dual-ground-truth evaluation to separate
statistical from causal validation. We perform a head-level analysis to
interpret recovered edges against the known GPT-2 induction circuit
\citep{olsson2022induction}. Our key findings are: (1) GES-BIC achieves the
best precision--speed trade-off among PGM methods; (2) statistical and
ablation methods find fundamentally different edges (Jaccard $\approx 0.10$);
and (3) feature selection determines which circuits are discoverable, with the
5\%--95\% firing-rate filter recovering a broad semantic circuit rather than
the induction circuit.

% ─────────────────────────────────────────────────────────────────────────────
\section{Background and Related Work}
% ─────────────────────────────────────────────────────────────────────────────

\subsection{Sparse Autoencoders}

Raw neurons in transformer models are \emph{polysemantic}: a single neuron
responds to many unrelated concepts due to superposition
\citep{elhage2022superposition}. Sparse Autoencoders (SAEs) learn an
overcomplete dictionary of approximately monosemantic features. Given a residual
stream vector $\mathbf{x} \in \mathbb{R}^{d}$, the SAE encoder computes:
\[
  \mathbf{f}(\mathbf{x}) = \text{ReLU}(W_{\text{enc}}\,\mathbf{x} + \mathbf{b})
  \in \mathbb{R}^{d_{\text{sae}}}
\]
where $d_{\text{sae}} \gg d$ ($d{=}768$, $d_{\text{sae}}{=}24{,}576$).
The reconstruction $\hat{\mathbf{x}} = W_{\text{dec}}\,\mathbf{f}(\mathbf{x}) + \mathbf{b}_{\text{dec}}$
is trained to be accurate while keeping $\mathbf{f}$ sparse, encouraging each
feature to encode one concept. We use the \texttt{gpt2-small-res-jb} release
from \texttt{sae\_lens} \citep{bloom2024saelens}, hooking at
\texttt{blocks.1.hook\_resid\_pre} (L1) and \texttt{blocks.2.hook\_resid\_pre} (L2).

\subsection{The GPT-2 Induction Circuit}

\citet{olsson2022induction} identified an induction circuit enabling
in-context learning. On sequences with repeated pattern $A \ldots B \ldots A$,
the model predicts $B$ via: \textbf{(1)} prev-token heads (L0H0, L0H1) copy
the previous token's identity into the current residual stream; \textbf{(2)}
induction heads (L1H5, L1H6) use this information to predict the token that
followed $A$ previously. This circuit operates precisely between the L1 and
L2 residual streams, making it an ideal ground truth for our task.

\subsection{Related Work on Automated Circuit Discovery}

\citet{conmy2023acdc} introduced Automated Circuit DisCovery (ACDC), which
iteratively ablates attention edges while monitoring a task-specific metric to
prune a full transformer computational graph. This approach requires a
pre-specified task and metric. \citet{marks2025sparse} recover sparse feature
circuits via integrated-gradient attribution through the SAE feature graph,
producing interpretable edge weights. Both methods treat circuit discovery as
a \emph{task-specific} problem. Our approach is complementary: we use general
PGM structure learning on activation data without specifying a task,
investigating whether the circuits implicit in the data's statistical structure
can be recovered without supervision.

% ─────────────────────────────────────────────────────────────────────────────
\section{Problem Formulation}
% ─────────────────────────────────────────────────────────────────────────────

Let $\mathbf{X} \in \mathbb{R}^{N \times p}$ with $N{=}5{,}000$ and $p{=}76$
($p_1{=}30$ L1 features, $p_2{=}46$ L2 features). We partition features as
$\mathcal{V} = \mathcal{L}_1 \cup \mathcal{L}_2$ and seek a directed graph
$\mathcal{G} = (\mathcal{V}, \mathcal{E})$.

We make two structural assumptions. First, we assume the joint distribution of
$\mathbf{X}$ is \emph{Markov} with respect to some DAG over $\mathcal{V}$
---the standard assumption underlying both PC and GES. Second, we use the known
\emph{architectural causal ordering}: layer 1 computation causally precedes
layer 2, so we restrict $\mathcal{E} \subseteq \mathcal{L}_1 \times
\mathcal{L}_2$. An edge $(i, j) \in \mathcal{E}$ means L1 feature $i$ directly
influences L2 feature $j$ through the layer-1 attention computation.

% ─────────────────────────────────────────────────────────────────────────────
\section{Methods}
% ─────────────────────────────────────────────────────────────────────────────

\subsection{PC Algorithm with Architectural Orientation}

The PC algorithm \citep{spirtes2000causation} learns a skeleton via conditional
independence testing, then orients edges via v-structure detection and Meek's
rules \citep{meek1995causal}. We run the skeleton phase with two tests, setting
$|S| \leq 3$ to limit exponential growth in conditioning-set enumeration
(depth-3 bounds the worst-case complexity to $\mathcal{O}(p^5)$).

\paragraph{PC-FisherZ.}
For continuous activations, we test partial correlation $\hat{\rho}_{ij \mid S}$
obtained from the precision submatrix $\hat{\Omega} = \hat{\Sigma}^{-1}_{[i,j]\cup S}$:
\[
  \hat{\rho}_{ij \mid S} = \frac{-\hat{\Omega}_{ij}}{\sqrt{\hat{\Omega}_{ii}\hat{\Omega}_{jj}}}
\]
The Fisher-Z statistic $z = \frac{1}{2}\ln\frac{1+\hat{\rho}}{1-\hat{\rho}}$ is
asymptotically $\mathcal{N}(0,\, (n{-}|S|{-}3)^{-1})$ under independence.
Edge $(i,j)$ is removed when the two-tailed p-value exceeds $\alpha$.
\emph{Note:} SAE activations are ReLU outputs and thus zero-inflated and
non-Gaussian. Fisher-Z is asymptotically valid but may lose power at small
sample sizes relative to the true conditional distribution.

\paragraph{PC-GSq.}
For binary activations (thresholded at zero), we use the G-test accumulated
over strata defined by the conditioning set:
\[
  G = 2 \sum_{\text{strata}} \sum_{a,b \in \{0,1\}} O_{ab} \ln\!\frac{O_{ab}}{E_{ab}}
\]
under a $\chi^2$ null. This makes no distributional assumption but discards
activation magnitude information.

\paragraph{Orientation.}
Rather than running v-structure detection and Meek's rules (which assume an
unknown causal ordering), we orient all cross-layer skeleton edges $i \to j$
($i \in \mathcal{L}_1$, $j \in \mathcal{L}_2$) using the known architectural
ordering. This is more informative than Meek's rules and directly exploits
the transformer's feed-forward structure.

\subsection{GES-BIC (Custom Fast Implementation)}

Greedy Equivalence Search \citep{chickering2002ges} optimises the BIC score
over equivalence classes of DAGs in two phases. We score node $j$ given parent
set $\text{Pa}(j)$ as:
\[
  \text{BIC}(j, \text{Pa}(j)) = -n \ln\!\!\left(\frac{\text{RSS}(j,\,\text{Pa}(j))}{n}\right)
  - |\text{Pa}(j)|\cdot\ln n
\]
where $\text{RSS}$ is the OLS residual sum of squares.

\paragraph{Schur-complement acceleration.}
Precomputing $\mathbf{C} = \mathbf{X}_c^\top \mathbf{X}_c$ (mean-centered)
once enables a closed-form RSS:
\[
  \text{RSS}(j, \text{Pa}) = C_{jj} -
  \mathbf{c}_{\text{Pa},j}^\top \mathbf{C}_{\text{Pa},\text{Pa}}^{-1}\,
  \mathbf{c}_{\text{Pa},j}
\]
All score evaluations use a single linear solve on $\mathbf{C}$, eliminating
repeated data access. This reduces per-iteration cost from
$\mathcal{O}(Np^2)$ to $\mathcal{O}(p^3)$ (dominated by the solve),
yielding \textbf{36s} total runtime for $p{=}76$, $N{=}5{,}000$---versus hours
for causal-learn's GES implementation. The forward phase adds 819 edges; the
backward phase prunes to 157 final cross-layer edges.

\subsection{Mean Ablation Patching}

As a causal proxy baseline, we estimate the influence of L1 feature $i$ on
L2 feature $j$ by replacing feature $i$'s activation with its dataset mean
$\bar{x}_i$ and measuring the downstream change:
\[
  A_{ij} = \frac{1}{N}\sum_{n=1}^{N}\!\left|
    f_j\!\left(\mathbf{x}_n;\,\tilde{\mathbf{x}}^{(i)}_n\right)
    - f_j(\mathbf{x}_n)\right|
\]
Edges are thresholded at the 95th percentile of $\{A_{ij}\}$, yielding 69
edges. This is \emph{mean-ablation sensitivity}---a linear, first-order causal
proxy---and is distinct from the integrated-gradient attribution of
\citet{marks2025sparse}, which computes gradient $\times$ activation through
the full SAE computational graph.

% ─────────────────────────────────────────────────────────────────────────────
\section{Experimental Setup}
% ─────────────────────────────────────────────────────────────────────────────

\paragraph{Data.}
We pass 5,000 uniformly random token sequences of length 40 through GPT-2
Small (\texttt{HookedTransformer} \citep{nanda2022transformerlens}, no
preprocessing). The random-prompt design avoids task-specific bias, at the
cost of not specifically activating circuits such as induction that require
structured repetition. We encode the last-token residual stream at both SAE
hooks, producing activations $\mathbf{X} \in \mathbb{R}^{5000 \times 49152}$
per layer before filtering.

\paragraph{Feature selection.}
We retain features with firing rate (fraction of tokens with activation $> 0$)
in $[5\%, 95\%]$, removing always-on and always-off features. This yields 30
L1 and 46 L2 features ($p = 76$ total). An important scope limitation:
circuit-specific features that fire rarely (e.g., induction features requiring
$A \ldots B \ldots A$ patterns) are systematically excluded by this filter.

\paragraph{Validation.}
We stratify 40 edges for ground-truth validation: 20 PGM-only (in any PGM
method but not ablation), 10 ablation-only, and 10 from the agreement set.
This stratification ensures coverage across method disagreement regions.
For each edge we apply two independent criteria.

\textbf{Interventional GT}: zero-ablate L1 feature $i$ on 500 held-out
prompts; confirm if $\langle|\Delta f_j|\rangle > 0.10 \cdot \langle
f_j\rangle_{\text{clean}}$.

\textbf{Observational GT}: compute Spearman rank correlation $\rho_s(X_i,
X_j)$ over 5,000 prompts; confirm if $|\rho_s| > 0.20$ and $p < 0.01$.

\emph{Intervention Agreement} (IA) under ground truth $\mathcal{T}$ is:
\[
  \text{IA}(\mathcal{M}, \mathcal{T}) =
  \frac{|\{e \in \hat{\mathcal{E}}_\mathcal{M} \cap
  \mathcal{E}_{\text{val}} : \mathcal{T}(e) = 1\}|}
  {|\hat{\mathcal{E}}_\mathcal{M} \cap \mathcal{E}_{\text{val}}|}
\]
With only 40 validated edges, IA estimates carry confidence intervals of
$\approx \pm 20$ percentage points (95\% CI via binomial proportion); method
comparisons where $N_{\text{val}} < 20$ should not be over-interpreted.

% ─────────────────────────────────────────────────────────────────────────────
\section{Results}
% ─────────────────────────────────────────────────────────────────────────────

\subsection{Edge Counts and Intervention Agreement}

\begin{table}[t]
\centering
\caption{Method comparison under interventional (ablation) ground truth.
  $N_{\text{val}}$: number of predicted edges that fell in the validated
  sample. $\dagger$: Abl-IA for Mean Ablation is upward-biased; see
  Section~\ref{sec:bias}. All IA estimates carry $\approx\pm 20$pp 95\% CI.}
\label{tab:results}
\small
\setlength{\tabcolsep}{5pt}
\begin{tabular}{lrrrr}
\toprule
Method & \# Edges & Abl-IA & $N_{\text{val}}$ & Runtime \\
\midrule
PC-FisherZ $\alpha{=}0.01$ & 184 & 0.667 & 18 & 33s \\
PC-FisherZ $\alpha{=}0.05$ & 223 & 0.619 & 21 & 64s \\
PC-FisherZ $\alpha{=}0.10$ & 255 & 0.619 & 21 & 98s \\
PC-GSq $\alpha{=}0.01$     & 160 & 0.643 & 14 & 247s \\
PC-GSq $\alpha{=}0.05$     & 200 & 0.611 & 18 & 499s \\
PC-GSq $\alpha{=}0.10$     & 226 & 0.600 & 20 & 682s \\
\textbf{GES-BIC}           & \textbf{157} & \textbf{0.667} & 15 & \textbf{36s} \\
Mean Ablation 95th         & 69  & 0.850$^\dagger$ & 20 & --- \\
\bottomrule
\end{tabular}
\end{table}

\paragraph{GES-BIC achieves the best precision--speed trade-off.}
GES-BIC recovers the fewest cross-layer edges (157) while matching the best
PC-FisherZ Ablation IA (0.667), in 36 seconds. Equal IA with fewer edges
implies higher precision: GES-BIC produces less noise per confirmed edge than
any PC variant. PC edge counts increase monotonically with $\alpha$ while Abl-IA
decreases, confirming that looser testing accumulates unconfirmed false positives.

\paragraph{PC-FisherZ vs.\ PC-GSq.}
FisherZ finds slightly more edges at significantly lower runtime. GSq's binary
representation discards magnitude information; while this removes the
Gaussianity assumption, it reduces statistical power without an apparent
precision gain. The IA gap (FisherZ 0.619, GSq 0.611 at $\alpha{=}0.05$)
is within the CI but directionally consistent.

\paragraph{Absolute IA in context.}
A naive random baseline---selecting edges uniformly at random from the 1380
possible cross-layer pairs---would yield an Ablation IA equal to the proportion
of all true edges among all possible edges, which is unknown but expected to
be substantially below 0.5 given the sparse nature of neural circuits.
The PGM methods achieving 0.61--0.67 thus represent genuine signal, despite
the wide confidence intervals.

\subsection{Pairwise Method Agreement}

\begin{table}[t]
\centering
\caption{Pairwise Jaccard similarity $|A \cap B| / |A \cup B|$.}
\label{tab:jaccard}
\small
\begin{tabular}{lrrrr}
\toprule
 & PC-FZ & PC-GSq & GES-BIC & Mean Abl \\
\midrule
PC-FZ    & 1.00 & 0.43 & 0.26 & 0.11 \\
PC-GSq   & 0.43 & 1.00 & 0.19 & 0.10 \\
GES-BIC  & 0.26 & 0.19 & 1.00 & 0.11 \\
Mean Abl & 0.11 & 0.10 & 0.11 & 1.00 \\
\bottomrule
\end{tabular}
\end{table}

The Jaccard matrix reveals three tiers of agreement. (1)~\emph{Within-paradigm}:
PC-FZ and PC-GSq share 0.43---same skeleton algorithm, different test, moderate
overlap. (2)~\emph{Cross-paradigm statistical}: PC vs.\ GES share 0.19--0.26---
both statistical, but GES's global BIC criterion is more conservative than
per-edge hypothesis tests, yielding a more parsimonious skeleton. (3)~\emph{
Statistical vs.\ causal}: PGM methods vs.\ Mean Ablation share only 0.10.

This low overlap is not noise---it reflects a fundamental conceptual
distinction. Statistical methods identify features that \emph{co-activate}
across prompts, including both direct causes and features driven by a common
upstream cause. Ablation methods identify features where one \emph{directly
drives} the other when intervened upon. These are different causal relations,
and the data shows they select very different edge sets.

\subsection{Evaluation Bias}
\label{sec:bias}

The 0.850 Ablation IA for Mean Ablation vs.\ 0.60--0.67 for PGM methods might
suggest ablation-based discovery is superior. This comparison is structurally
unfair: both zero-ablation validation and Mean Ablation use activation
disruption as their lens. An edge $(i,j)$ found because ablating $i$ changes
$j$'s mean activation will trivially also be confirmed by zero-ablation
validation. PGM methods, conversely, flag statistical co-occurrences including
shared-cause correlations; ablation validation correctly rejects these, since
zeroing one branch of a shared-cause pair does not affect the other. The
upward bias on Mean Ablation's Abl-IA is therefore structural, not inferential.

The Spearman Observational IA addresses this: Spearman correlation shares no
inductive bias with any method (it is purely observational), providing a fair
comparison across both statistical and causal methods. We describe this
protocol in Section~5 and defer numerical results to a longer version of this
work, where a full validation run is completed on all methods simultaneously.

% ─────────────────────────────────────────────────────────────────────────────
\section{Head-Level Circuit Recovery}
% ─────────────────────────────────────────────────────────────────────────────

\subsection{Feature-to-Head Attribution}

We attribute each SAE feature to its dominant attention head by projecting the
feature's encoder direction onto head-output space via $W_O$:
\[
  \alpha_{h,f}^{(L)} = \frac{1}{N}\sum_{n=1}^{N}
  \mathbf{z}_{n,h}^{(L)}\,W_O^{(L)}[h]\cdot\mathbf{w}_{\text{enc},f}^{(L)}
\]
where $\mathbf{z}_{n,h}^{(L)}$ is the pre-$W_O$ output of head $h$ on prompt
$n$. The dominant head assignment is $h^*(f) = \arg\max_h |\alpha_{h,f}|$.
Feature-level edges $(i,j)$ are then aggregated to head-level counts:
$M[h^*(i), h^*(j{-}p_1)] {+}{=} 1$.

\subsection{Results and Interpretation}

Head assignments are highly concentrated: 20 of 30 L1 features map to L0H9;
20 of 46 L2 features map to L1H2 and 20 to L1H10. The known induction circuit
(L0H0/H1 $\to$ L1H5/H6) scores an induction fraction of \textbf{0.000} across
all methods.

\paragraph{Why induction was not recovered.}
The 5\%--95\% filter selects broadly-active features. The induction circuit
activates on rare $A\ldots B\ldots A$ patterns; induction-specific SAE features
have low firing rates and are excluded. This is a scope limitation of feature
selection, not of the structure-learning algorithms: given induction-specific
features as input, the algorithms would have the opportunity to recover the
circuit.

\paragraph{Consistent alternate circuit.}
All four methods agree on L0H9 $\to$ L1H2/H10. L0H9 in GPT-2 Small is
associated with broad positional and syntactic processing (carrying
token-position signals across all inputs), consistent with its domination
of broadly-active L1 features. L1H2 and L1H10 are known to process semantic
and token-type information at layer 1. The recovered circuit plausibly
reflects broad semantic content routing from layer 0 into layer 1
representations---a real computational pattern, distinct from but not less
meaningful than the induction mechanism.

The cross-method consistency of this finding is a positive result: four
independent algorithms with different assumptions all identify the same
head-level structure, suggesting genuine signal rather than method-specific
artefacts.

% ─────────────────────────────────────────────────────────────────────────────
\section{Discussion}
% ─────────────────────────────────────────────────────────────────────────────

\paragraph{Scalability constraints motivate targeted feature selection.}
The principled extension of our work is to relax the firing-rate filter
to include induction-specific features. Table~\ref{tab:complexity} shows this
is combinatorially infeasible without algorithmic changes. The solution is
\emph{targeted feature selection}: compute head--feature attribution
$\alpha_{h,f}$ for all 24,576 SAE features (a single forward pass), select
the top-$k$ features per known circuit head (e.g.\ L0H0/H1, L1H5/H6), and
run GES on the resulting $\approx 80$-node graph. This keeps runtime at $\sim$36s
while directly targeting induction-relevant features.

\begin{table}[h]
\centering
\caption{Estimated runtimes as firing-rate filter is relaxed.}
\label{tab:complexity}
\small
\begin{tabular}{llll}
\toprule
$p$ & Filter & PC-FisherZ & GES-BIC \\
\midrule
76 (current) & 5--95\%    & 64s      & 36s \\
$\sim$300    & 1--99\%    & $\sim$8h & $\sim$15min \\
$\sim$800    & 0.5--99.5\%& $\sim$100 days & $\sim$5h \\
49,152 (all) & none       & centuries & weeks \\
\bottomrule
\end{tabular}
\end{table}

\paragraph{PGM structure learning is viable for circuit discovery.}
All three PGM methods produce partially-confirmed edge sets without any
task-specific supervision. GES-BIC's combination of global score optimisation
and Schur-complement acceleration makes it practically competitive with the
causal proxy baseline on Ablation IA while maintaining methodological rigour.

\paragraph{Statistical vs.\ causal methods find different circuits.}
The Jaccard $\approx 0.10$ between PGM and Mean Ablation is a substantive
finding. It rules out the hypothesis that statistical co-occurrence
straightforwardly proxies causal influence in transformer activations. Both
types of edges are scientifically meaningful---co-activation structure reflects
shared computational roles, while causal edges reflect direct information
routing---but they should not be conflated.

\paragraph{Limitations.}
(1) PC runs the skeleton phase only; within-layer edge orientation would
require v-structure detection and Meek's rules.
(2) With 40 validated edges, all IA estimates carry $\approx\pm 20$pp 95\% CIs;
method differences are suggestive, not conclusive.
(3) Fisher-Z assumes Gaussianity; SAE activations are ReLU-gated and
zero-inflated, which may inflate type-I error at small sample sizes.
(4) Mean Ablation captures only first-order effects; nonlinear or multi-hop
causal paths are not measured.
(5) All experiments use GPT-2 Small on random prompts; generalisation to
other models or task-specific distributions is not assessed.

% ─────────────────────────────────────────────────────────────────────────────
\section{Conclusion}
% ─────────────────────────────────────────────────────────────────────────────

We have demonstrated that classical PGM structure learning---PC and GES---can
recover partially causally-validated feature circuits in a language model from
activation data alone, without task-specific supervision. GES-BIC is the
recommended method: it achieves the best precision--speed trade-off, avoids
the false-positive accumulation of per-edge hypothesis testing, and scales
gracefully via Schur-complement score evaluation.

The low Jaccard similarity between statistical and ablation methods
($J \approx 0.10$) is a key negative result: co-activation structure does not
straightforwardly proxy direct causal influence. We argue this motivates
a \emph{two-layer} evaluation protocol---interventional for causal edges,
observational for statistical dependencies---and that single-metric evaluation
conflating the two produces misleading conclusions.

The failure to recover the induction circuit reveals that feature selection is
the binding constraint, not the structure-learning algorithms. Targeted
head-attribution-seeded feature selection is the natural next step: it would
preserve the computational efficiency of our approach while directly targeting
circuit-specific features, and represents a promising direction toward scalable,
automated, supervision-free circuit discovery in large language models.

% ─────────────────────────────────────────────────────────────────────────────
\bibliographystyle{plainnat}
\begin{thebibliography}{99}

\bibitem[Bloom et al.(2024)]{bloom2024saelens}
Joseph Bloom, David Chanin, Curt Tigges, and Neel Nanda.
\newblock \texttt{SAELens}.
\newblock \emph{GitHub}, 2024.
\newblock \url{https://github.com/jbloomaus/SAELens}.

\bibitem[Chickering(2002)]{chickering2002ges}
David Maxwell Chickering.
\newblock Optimal structure identification with greedy search.
\newblock \emph{JMLR}, 3:507--554, 2002.

\bibitem[Conmy et al.(2023)]{conmy2023acdc}
Arthur Conmy, Augustine N.\ Mavor-Parker, Aengus Lynch, Stefan Heimersheim,
and Adri\`{a} Garriga-Alonso.
\newblock Towards automated circuit discovery for mechanistic interpretability.
\newblock In \emph{NeurIPS}, 2023.

\bibitem[Elhage et al.(2021)]{elhage2021framework}
Nelson Elhage, Neel Nanda, Catherine Olsson, et al.
\newblock A mathematical framework for transformer circuits.
\newblock \emph{Transformer Circuits Thread}, 2021.

\bibitem[Elhage et al.(2022)]{elhage2022superposition}
Nelson Elhage, Tristan Henighan, Nicholas Joseph, et al.
\newblock Toy models of superposition.
\newblock \emph{Transformer Circuits Thread}, 2022.

\bibitem[Marks et al.(2025)]{marks2025sparse}
Samuel Marks, Can Rager, Eric J.\ Michaud, Yonatan Belinkov, David Bau,
and Aaron Mueller.
\newblock Sparse feature circuits: Discovering and editing interpretable causal
graphs in language models.
\newblock \emph{ICLR}, 2025.

\bibitem[Meek(1995)]{meek1995causal}
Christopher Meek.
\newblock Causal inference and causal explanation with background knowledge.
\newblock In \emph{UAI}, pages 403--410, 1995.

\bibitem[Nanda and Lawrence(2022)]{nanda2022transformerlens}
Neel Nanda and Joseph Lawrence.
\newblock TransformerLens.
\newblock \emph{GitHub}, 2022.
\newblock \url{https://github.com/neelnanda-io/TransformerLens}.

\bibitem[Olah et al.(2020)]{olah2020zoom}
Chris Olah, Nick Cammarata, Ludwig Schubert, Gabriel Goh, Michael Petrov,
and Shan Carter.
\newblock Zoom in: An introduction to circuits.
\newblock \emph{Distill}, 2020.

\bibitem[Olsson et al.(2022)]{olsson2022induction}
Catherine Olsson, Nelson Elhage, Neel Nanda, Nicholas Joseph, Nova DasSarma,
Tom Henighan, Ben Mann, et al.
\newblock In-context learning and induction heads.
\newblock \emph{Transformer Circuits Thread}, 2022.

\bibitem[Radford et al.(2019)]{radford2019gpt2}
Alec Radford, Jeffrey Wu, Rewon Child, David Luan, Dario Amodei, and Ilya Sutskever.
\newblock Language models are unsupervised multitask learners.
\newblock \emph{OpenAI Blog}, 1(8):9, 2019.

\bibitem[Spirtes et al.(2000)]{spirtes2000causation}
Peter Spirtes, Clark Glymour, and Richard Scheines.
\newblock \emph{Causation, Prediction, and Search}.
\newblock MIT Press, 2nd ed., 2000.

\bibitem[Wang et al.(2022)]{wang2022ioi}
Kevin Wang, Alexandre Variengien, Arthur Conmy, Buck Shlegeris,
and Jacob Steinhardt.
\newblock Interpretability in the wild: A circuit for indirect object
identification in GPT-2 small.
\newblock \emph{ICLR}, 2023.

\end{thebibliography}

\end{document}
